\documentclass[../../../thesis.tex]{subfiles}
\begin{document}
\subsection{Anglický jazyk}
Šablóna FEIstyle podporuje slovenský a anglický jazyk.
Pre prácu v anglickom jazyku je potrebné túto skutočnosť nastaviť v preambule hlavného súboru \verb|thesis.tex| ako nepovinný parameter \verb|en| príkazu definície šablóny
\begin{trivlist}
\item \texttt{\textbackslash documentclass[bp,en]\{FEIstyle\}}
\end{trivlist}
Prvý parameter určuje, či ide o bakalársku (\verb|bp|) alebo diplomovú (\verb|dp|) prácu.

Anglická práca musí obsahovať po závere rezumé a na to je potrebné odstrániť komentár pred príkazom
\verb|\FEIresume{includes/resume}|.

Na jazykovú lokalizáciu používame balík \verb|babel|.
Ak sa v práci písanej v slovenčine nachádzajú výrazy v angličtine,
uvedieme takýto text do makra \verb|\foreignlanguage|, čím zabezpečíme správne medzery a~delenie slov.
Napríklad pri zavádzaní skratky AI môžeme napísať, že ide o anglický výraz pre umelú inteligencie \foreignlanguage{english}{Artificial Intelligence}.
Zapíšeme ho nasledujúcim spôsobom:
\begin{verbatim}
\foreignlanguage{english}{Artificial Intelligence}
\end{verbatim}

Ak nastavíme parametrom \verb|en| anglický jazyk ako hlavný, stane sa slovenčina cudzím jazykom.
\end{document}